\section{Introduction}
Trong bối cảnh điện khí hóa đang trở thành khuynh hướng tất yếu trong ngành giao thông vận tải và năng lượng tái tạo, nhu cầu về các dòng máy điện có mật độ công suất và hiệu suất vượt trội ngày càng trở nên cấp thiết . Trong đó, động cơ nam châm vĩnh cửu hướng trục (AFPM) đã khẳng định vị thế là một công nghệ then chốt nhờ những ưu thế đặc biệt về cấu trúc \cite{jenkins2023innovations,gieras2008axial,capponi2012recent,xu2021comparative}. 

Tuy nhiên, khác biệt hoàn toàn với các dòng máy điện hướng kính truyền thống, AFPM sở hữu đặc thù hình học với sự phân bổ từ trường biến thiên phi tuyến và phức tạp trong không gian ba chiều \cite{vatani2025axial,ni2013investigation}. Đặc tính này khiến các mô hình 2D đơn giản hóa không còn đủ khả năng đảm bảo độ tin cậy trong việc phân tích và đánh giá chính xác các đặc tính vận hành. 

Để đạt được độ chính xác cần thiết trong thực tế, các nhà thiết kế bắt buộc phải sử dụng đến phương pháp Phần tử hữu hạn 3D (3D-FEM). Mặc dù vậy, thách thức nảy sinh khi các quy trình thiết kế hiện đại thường đi kèm với bài toán tối ưu hóa đa mục tiêu sử dụng các thuật toán tìm kiếm toàn cục \cite{huner2023design,subotic2019weight,liu2025optimization}. Những thuật toán này đòi hỏi hàng ngàn chu kỳ đánh giá mô hình lặp lại, biến chi phí tài nguyên và thời gian thực thi cực lớn của 3D~FEM truyền thống trở thành một nút thắt cổ chai kỹ thuật đáng kể, cản trở tiến độ nghiên cứu và phát triển sản phẩm.


Các phương pháp tham số tập trung (Lumped Parameter Method - LPM), hay còn được biết đến với tên gọi mạng từ trở tương đương (Magnetic Equivalent Circuit - MEC), đã được nghiên cứu và phát triển từ rất sớm như một hướng tiếp cận kinh điển nhằm tối ưu hóa tốc độ tính toán thông qua việc mô hình hóa không gian bằng một hệ thống mạch từ cấu tạo từ các từ trở và nguồn sức từ động . Nhờ ưu thế vượt trội về mặt thời gian thực thi, phương pháp này đã được ứng dụng thành công và rộng rãi trên nhiều dòng máy điện từ truyền thống như động cơ nam châm vĩnh cửu bề mặt (SPMSM) \cite{manju2018magnetic,cui2025equivalent,zhang2019equivalent,yang2011surface}, nam châm chìm (IPMSM) \cite{kemmetmuller2014modeling,zhu2009analytical,tariq2010iron} cho đến các cấu trúc AFPM hiện đại \cite{zhang20213,6971543,6600909,7869440} . Tuy nhiên, các phương pháp này thường phụ thuộc đáng kể vào giả thiết hình học và quy luật dòng từ vốn dựa trên kinh nghiệm chủ quan của người thiết kế, dẫn đến những hạn chế khi áp dụng cho các cấu trúc có phân bổ từ trường phức tạp hoặc yêu cầu độ chính xác cục bộ cao.


Một nhánh nghiên cứu gần đây dựa trên việc rời rạc hóa không gian tính toán thành duy nhất một loại phần tử hình học thông qua lưới phi tham số (non-parametric mesh) đã mở ra hướng đi mới cho kỹ thuật mạng từ trở (MRN). Phương pháp này được giới thiệu ban đầu thông qua việc sử dụng các phần tử hình chữ nhật trong hệ tọa độ Cartesian để tự động hóa quá trình tạo lưới \cite{9473194}. Các công trình tiêu biểu tiếp theo đã phát triển kỹ thuật này bằng cách sử dụng các phần tử hình tam giác \cite{10018857,ouagued2016comparison}, hình quạt \cite{lu2023hybrid,10599821} cho phép mô hình hóa các dòng máy điện hướng kính một cách hiệu quả. Tuy nhiên, các nghiên cứu này chỉ dừng lại ở miền tính toán hai chiều (2D), khiến việc khai thác kỹ thuật này trong không gian ba chiều còn nhiều hạn chế.

Xuất phát từ bối cảnh đó, bài báo này đề xuất phương pháp mạng từ trở 3D phân tích nút lưới phi tham số (Non-Parametric 3D-Nodal MRN) cho các máy điện quay tổng quát. Phương pháp đề xuất, gọi tắt là 3D MRN, thực hiện rời rạc hóa miền tính toán thành các phần tử $E$ trong không gian 3D tích hợp đầy đủ đặc tính vật liệu và nguồn kích thích. Bằng việc kết hợp kỹ thuật phân tích nút với thuật toán tìm kiếm giảm chấn (damping search algorithm), 3D MRN đảm bảo tính hội tụ và ổn định cho các hệ thống quy mô lớn. Giải pháp này cho phép mô phỏng chính xác các đặc tính điện từ 3D với tốc độ tính toán nhanh, hỗ trợ đắc lực cho quy trình thiết kế và tối ưu hóa máy điện.

Bài báo được trình bày thành 6 phần: Phần I giới thiệu tổng quan; Phần II trình bày phương pháp xây dựng mô hình 3D MRN và chiến lược rời rạc hóa miền tính toán bằng phần tử $E$; Phần III mô tả các kỹ thuật hậu xử lý; Phần IV thực hiện xác minh độ chính xác qua so sánh với 3D~FEM và thực nghiệm; Phần V thảo luận về các kỹ thuật tối ưu hóa bộ giải; và Phần VI đưa ra các kết luận của nghiên cứu.